\chapter{Fixed-Root Exclusion from Quantitative Kernel Curvature}
\label{chap:fixed-root-exclusion-kernel-curvature}

\status{Exact analytic exclusion proved: $F(-1/4)>0$. Hence the negative-inversion paired quotient spectrum has no fixed roots, $|P_J|$ is even, and $|P_N|$ is divisible by four. SOH-G003 and RH remain open.}

Chapters~\ref{chap:quotient-negative-inversion-zero-set-no-go} and
\ref{chap:paired-spectrum-quotient-correspondence} reduced the only possible
fixed exceptional quotient root to
\[
 w=-\frac14.
\]
The positive fixed value $+1/4$ had already been excluded by positivity of the
coefficients of $F$.  This chapter removes the remaining negative fixed value
analytically, without using a zero table.

\section{Attribution boundary}

Gershon proved strict log-concavity of the classical Riemann Xi kernel
$\Phi$ independently in 2026~\cite{gershon2026logconcavity}.  The present
argument does not claim priority for that fact.  What is used here is the
explicit quantitative SOH-G004 channel/mixture estimate from
Chapter~\ref{ch:global-log-concavity}: each theta channel satisfies
\[
 g_n''<-12,
\]
while the mixture slope variance obeys the conservative bound
\[
 \operatorname{Var}_p(g_n')<2.
\]
Those retained constants yield a numerical curvature margin stronger than
mere strict log-concavity and are the input to the fixed-root exclusion below.

\section{Modular evenness of the kernel}

Let
\[
 \Theta(x)=\sum_{n\in\mathbb Z}e^{-\pi n^2x},
 \qquad x>0.
\]
Jacobi modularity gives
\begin{equation}
 \Theta(x)=x^{-1/2}\Theta(1/x).
 \label{eq:g017-theta-modular}
\end{equation}
Writing $x=e^{2y}$, the positive Xi kernel from
Chapter~\ref{chap:positive-riemann-kernel} is equivalently
\begin{equation}
 \Phi(y)=4x^{9/4}\Theta''(x)+6x^{5/4}\Theta'(x).
 \label{eq:g017-kernel-theta}
\end{equation}

Put $t=1/x$.  Differentiating \cref{eq:g017-theta-modular} once and twice
gives
\begin{align}
 \Theta'(t)
 &=-x^{5/2}\Theta'(x)-\frac12x^{3/2}\Theta(x),
 \label{eq:g017-theta-first}\
 \Theta''(t)
 &=x^{9/2}\Theta''(x)
   +3x^{7/2}\Theta'(x)
   +\frac34x^{5/2}\Theta(x).
 \label{eq:g017-theta-second}
\end{align}
Substituting these identities into the right-hand side of
\cref{eq:g017-kernel-theta} at $t=1/x$ yields
\begin{align*}
4t^{9/4}\Theta''(t)+6t^{5/4}\Theta'(t)
&=4x^{9/4}\Theta''(x)+6x^{5/4}\Theta'(x).
\end{align*}
The terms proportional to $\Theta(x)$ cancel exactly. Therefore
\begin{equation}
 \boxed{\Phi(-y)=\Phi(y).}
 \label{eq:g017-kernel-even}
\end{equation}
In particular,
\begin{equation}
 \boxed{\Phi'(0)=0.}
 \label{eq:g017-kernel-centre-derivative}
\end{equation}

\section{Quantitative strong log-concavity}

The exact mixture identity from Chapter~\ref{ch:global-log-concavity} is
\[
 (\log\Phi)''
 =\sum_n p_ng_n''+\operatorname{Var}_p(g_n').
\]
Combining
\[
 g_n''<-12
\]
with
\[
 \operatorname{Var}_p(g_n')<2
\]
gives the strict global estimate
\begin{equation}
 \boxed{(\log\Phi)''<-10.}
 \label{eq:g017-strong-log-concavity}
\end{equation}

Set
\[
 V(y)=-\log\Phi(y).
\]
Then $V''(y)>10$.  Equation~\eqref{eq:g017-kernel-centre-derivative}
gives $V'(0)=0$, hence for every $y>0$,
\begin{equation}
 \boxed{V'(y)>10y.}
 \label{eq:g017-potential-slope}
\end{equation}

\section{A direct second-moment bound}

Define
\[
 m_0=\int_0^\infty\Phi(y)\,dy,
 \qquad
 m_2=\int_0^\infty y^2\Phi(y)\,dy.
\]
The explicit theta-kernel formula decays super-exponentially, so
$y\Phi(y)\to0$ as $y\to\infty$. Since $V'\Phi=-\Phi'$, the slope bound
\cref{eq:g017-potential-slope} gives
\begin{align}
10m_2
&<\int_0^\infty yV'(y)\Phi(y)\,dy\\
&=-\int_0^\infty y\Phi'(y)\,dy\\
&=m_0.
\end{align}
Thus
\begin{equation}
 \boxed{m_2<\frac{m_0}{10}.}
 \label{eq:g017-moment-bound}
\end{equation}
No general concentration theorem is required; the estimate follows directly
from the quantitative logarithmic curvature and one integration by parts.

\section{Exclusion of the negative fixed quotient value}

For $w=-1/4$ choose $z=i/2$.  The positive-kernel representation gives
\begin{equation}
 F(-1/4)
 =\xi\!\left(\frac12+\frac{i}{2}\right)
 =\int_0^\infty\Phi(y)\cos(y/2)\,dy.
 \label{eq:g017-fixed-value-integral}
\end{equation}
The elementary global inequality
\[
 \cos x\ge1-\frac{x^2}{2}
\]
implies
\[
 \cos(y/2)\ge1-\frac{y^2}{8}.
\]
Therefore, using \cref{eq:g017-moment-bound},
\begin{align}
 F(-1/4)
 &\ge m_0-\frac{m_2}{8}\\
 &>m_0\left(1-\frac1{80}\right).
\end{align}
Since $m_0=F(0)=\xi(1/2)>0$ by the positive-kernel theorem,
\begin{theorem}[SOH-G017 fixed-root exclusion]
\label{thm:g017-fixed-root-exclusion}
\begin{equation}
 \boxed{
 F(-1/4)>\frac{79}{80}F(0)>0.
 }
 \label{eq:g017-main-lower-bound}
\end{equation}
In particular,
\[
 \boxed{F(-1/4)\ne0.}
\]
\end{theorem}

\section{Exceptional-orbit consequences}

The G015 involution
\[
 J(w)=\frac1{16w}
\]
has exactly two fixed values, $w=\pm1/4$.  Coefficient positivity already
gives $F(1/4)>0$, and \cref{thm:g017-fixed-root-exclusion} now gives
$F(-1/4)>0$.  Hence the finite paired-root set $P_J$ has no fixed points.
It is therefore a disjoint union of non-fixed two-cycles, so
\begin{equation}
 \boxed{|P_J|\equiv0\pmod2.}
 \label{eq:g017-PJ-even}
\end{equation}
G016 proved $|P_N|=2|P_J|$, and therefore
\begin{equation}
 \boxed{|P_N|\equiv0\pmod4.}
 \label{eq:g017-PN-mod4}
\end{equation}

The G012 geometric fixed pair
\[
 \left\{\frac12+\frac i2,\frac12-\frac i2\right\}
\]
is the fiber over $w=-1/4$.  It follows immediately that neither point is a
zero of $\xi$.

\section{Proof firewall}

SOH-G017 proves:
\begin{itemize}
 \item the modular evenness $\Phi(-y)=\Phi(y)$ and $\Phi'(0)=0$;
 \item the quantitative consequence $(\log\Phi)''<-10$ from the existing
       G004 constants;
 \item $m_2<m_0/10$;
 \item $F(-1/4)>(79/80)F(0)>0$;
 \item the paired quotient spectrum has no fixed roots;
 \item $|P_J|$ is even and $|P_N|$ is divisible by four;
 \item the G012 fixed pair $1/2\pm i/2$ is not in the xi zero set.
\end{itemize}

It does not prove:
\begin{itemize}
 \item that $P_J$ or $P_N$ is empty;
 \item that non-fixed paired two-cycles do not exist;
 \item real-rootedness of $F$;
 \item PF$_\infty$;
 \item SOH-G003;
 \item the Riemann Hypothesis.
\end{itemize}
